\documentclass[a4paper,11pt]{article}
\usepackage[T1]{fontenc}
\usepackage[utf8]{inputenc}
\usepackage[polish]{babel}
\usepackage{lmodern}
\usepackage{amsmath,amssymb}
\usepackage[margin=2.2cm]{geometry}
\pagestyle{plain}
\begin{document}
\begin{center}
{\Large\bfseries Przykładowy arkusz egzaminacyjny z matematyki}\\[4pt]
Poziom podstawowy -- materiał ćwiczeniowy (plik demonstracyjny)
\end{center}
\noindent\textbf{Instrukcja dla zdającego.} Sprawdź, czy arkusz zawiera wszystkie strony.
Rozwiązania zapisz czytelnie w miejscu na to przeznaczonym. Pamiętaj: zażółć gęślą jaźń --
ZAŻÓŁĆ GĘŚLĄ JAŹŃ.

\section*{Zadanie 1. (0--1)}
Liczba $\sqrt[3]{54}-\sqrt[3]{2}$ jest równa
\[
A.\ \sqrt[3]{52}\qquad B.\ 3\qquad C.\ 2\sqrt[3]{2}\qquad D.\ 2
\]
\section*{Zadanie 2. (0--2)}
Rozwiąż nierówność $2x^2-7x+3\geq 0$. Wynik zapisz w postaci sumy przedziałów.
\[
\Delta=b^2-4ac=49-24=25,\qquad x_{1,2}=\frac{7\pm 5}{4}
\]
\section*{Zadanie 3. (0--1)}
Ciąg arytmetyczny $(a_n)$ określony jest wzorem $a_n=3n-2$ dla $n\geq1$.
Suma $\displaystyle\sum_{k=1}^{10} a_k$ jest równa $\_\_\_\_$.
\section*{Zadanie 4. (0--3)}
Funkcja $f(x)=\dfrac{x^2-4}{x+2}$ jest określona dla $x\in\mathbb{R}\setminus\{-2\}$.
Wykaż, że $f(x)=x-2$ i oblicz $\int_0^1 f(x)\,dx$.
\vfill
\noindent Miejsce na obliczenia (nie podlega ocenie): \hrulefill
\newpage
\section*{Zadanie 5. (0--2)}
Kąt $\alpha\in(0^\circ,90^\circ)$ spełnia warunek $\sin\alpha=\tfrac{3}{5}$. Oblicz $\operatorname{tg}\alpha$.
\section*{Zadanie 6. (0--4)}
W trójkącie prostokątnym przyprostokątne mają długości $a=6$ i $b=8$.
Oblicz długość przeciwprostokątnej $c=\sqrt{a^2+b^2}$ oraz pole koła wpisanego w ten trójkąt.

\vspace{2cm}
\noindent Odpowiedź: \hrulefill
\end{document}
